\lyseis{Α3.2}{Απόλυτη τιμή, σύγκριση, διάταξη, αντίθετοι}
\sol{1} α) $-\alpha$ \quad β) μεγαλύτερη \quad γ) μικρότερη \quad δ) αριθμητή \quad ε) μικρότερο.
\sol{2} Στη στήλη του $-3$: αντίθετος $3$, απόλυτη τιμή $3$. Για τον $0{,}35$: $-0{,}35$ και $0{,}35$. Για τον $\frac{1}{9}$: $-\frac{1}{9}$ και $\frac{1}{9}$. Στην τέταρτη στήλη ο αντίθετος είναι ο $-0{,}1$, άρα ο αριθμός είναι ο $0{,}1$ και η απόλυτη τιμή $0{,}1$. Στην πέμπτη ο αντίθετος είναι ο $\frac{7}{4}$, άρα ο αριθμός είναι ο $-\frac{7}{4}$ και η απόλυτη τιμή $\frac{7}{4}$. Για τον $-\frac{2}{5}$: αντίθετος $\frac{2}{5}$ και απόλυτη τιμή $\frac{2}{5}$.
\sol{3} α) $>$ (ομώνυμα, $21>12$). β) $<$ (ίδιος αριθμητής, μεγαλύτερος παρονομαστής, άρα μικρότερο κλάσμα). γ) $\frac{2}{3}=\frac{16}{24}<\frac{21}{24}=\frac{7}{8}$, άρα $<$. δ) $>$ (θετικός και αρνητικός). ε) $-0{,}8>-1$, άρα $>$. στ) $-\frac{8}{3}=-2{,}666\ldots>-3$, άρα $>$. ζ) $\frac{9}{8}=1{,}125<1{,}2$, άρα $<$.
\sol{4} α) $3{,}16<3{,}20$, άρα $<$. β) $0{,}092>0{,}087$ (στα εκατοστά $9>8$), άρα $>$. γ) $-1{,}36<-0{,}01$ (αρνητικοί, ο $-1{,}36$ έχει μεγαλύτερη απόλυτη τιμή), άρα $<$. δ) $<$ (αρνητικός και θετικός). ε) $-1{,}3=-1{,}30>-1{,}33$, άρα $>$. στ) $-1{,}10=-1{,}1$, άρα $=$.
\sol{5} $-3<-0{,}03<0<0{,}2<1<1{,}37$.
\sol{6} α) $-\frac{5}{4}<-\frac{3}{4}<-\frac{1}{4}<\frac{1}{4}<\frac{3}{4}<\frac{5}{4}$. β) Οι αρνητικοί: $-\frac{5}{2}=-2{,}5$, $-\frac{5}{3}=-1{,}666\ldots$, $-\frac{5}{7}=-0{,}714\ldots$ Οι θετικοί, με ίδιο αριθμητή: όσο μεγαλύτερος ο παρονομαστής, τόσο μικρότερο το κλάσμα: $\frac{5}{8}<\frac{5}{4}<\frac{5}{3}$. Άρα $-\frac{5}{2}<-\frac{5}{3}<-\frac{5}{7}<\frac{5}{8}<\frac{5}{4}<\frac{5}{3}$.
\sol{7} α) Μικρότερα του $1$ είναι τα $\frac{3}{16},\ \frac{7}{16},\ \frac{15}{16}$ και $\frac{3}{16}<\frac{7}{16}<\frac{15}{16}$. β) Μεγαλύτερα του $1$ είναι τα $\frac{9}{8}=1{,}125$, $\frac{5}{4}=1{,}25$, $\frac{9}{4}=2{,}25$ και $\frac{9}{8}<\frac{5}{4}<\frac{9}{4}$. γ) $\frac{3}{16}<\frac{7}{16}<\frac{15}{16}<\frac{9}{8}<\frac{5}{4}<\frac{9}{4}$.
\sol{8} α) Σωστό. β) Λάθος: $|-5|=5$. γ) Σωστό. δ) Λάθος: $-9<-2$. ε) Σωστό (ίδιος αριθμητής, μεγαλύτερος παρονομαστής, μικρότερο κλάσμα). στ) Σωστό.
\sol{9} Από την πιο κρύα: Φλώρινα $-10$, Μέτσοβο $-5{,}5$, Ιωάννινα $-3$, Θεσσαλονίκη $-1{,}5$, Αθήνα $7$.
\sol{10} α) Ύπνος $24:3=8$ ώρες, Σχολείο $24:4=6$ ώρες, Διάβασμα $24:8=3$ ώρες, Διασκέδαση $24:6=4$ ώρες. β) Ύπνος ($8$), Σχολείο ($6$), Διασκέδαση ($4$), Διάβασμα ($3$). γ) Μένουν $24-21=3$ ώρες, δηλαδή $\frac{3}{24}=\frac{1}{8}$ του εικοσιτετραώρου. Έλεγχος με κλάσματα: $\frac{8}{24}+\frac{6}{24}+\frac{3}{24}+\frac{4}{24}=\frac{21}{24}$.
\sol{11} Κέρδισε ο Βασίλης: $9{,}75<9{,}80$ και ο μικρότερος χρόνος είναι ο καλύτερος. Το λάθος του: συνέκρινε το $75$ με το $8$ σαν φυσικούς. Τα $9{,}8$ είναι $9{,}80$, δηλαδή $80$ εκατοστά, και $80>75$.
\sol{12} Το σχολείο είναι στο $0$. Το βιβλιοπωλείο ($-2$) απέχει $2$ μονάδες, δηλαδή $200$ μ. Ο φούρνος ($+3$) απέχει $300$ μ. Πιο κοντά είναι το βιβλιοπωλείο. Τα δύο μαγαζιά απέχουν $2+3=5$ μονάδες, δηλαδή $500$ μ.
\sol{13} $-11\,034<-430<0<2\,917$: τάφρος Μαριανών, Νεκρά Θάλασσα, επιφάνεια θάλασσας, Όλυμπος.
\sol{14} $\frac{2}{5}=0{,}4$ και $\frac{3}{5}=0{,}6$. Ανάμεσα βρίσκονται π.χ. οι $0{,}45$, $0{,}5=\frac{1}{2}$ και $0{,}55$. Υπάρχουν άπειροι άλλοι: δεν υπάρχει «επόμενος».
\sol{15} α) Π.χ. $-0{,}75$ και $-0{,}72$ (ο $-0{,}8<-0{,}75<-0{,}7$). β) $\frac{3}{10}=0{,}3$ και $\frac{2}{5}=0{,}4$: π.χ. $0{,}35=\frac{7}{20}$ και $0{,}33$. γ) $\frac{7}{9}=0{,}777\ldots$ και $\frac{5}{6}=0{,}8333\ldots$: π.χ. $0{,}8=\frac{4}{5}$ και $0{,}79$.
\sol{16} $|\alpha|=\alpha$ για κάθε θετικό και για το $0$. $|\alpha|=-\alpha$ για κάθε αρνητικό και για το $0$. Ο $\alpha$ και ο $-\alpha$ είναι συμμετρικοί ως προς το $0$, άρα απέχουν το ίδιο από το $0$ και $|\alpha|=|-\alpha|$.
\sol{17} α) $0$ και $1$. β) $-2$ και $-1$. γ) $\frac{7}{2}=3{,}5$: $3$ και $4$. δ) $-\frac{11}{3}=-3{,}666\ldots$: $-4$ και $-3$. ε) $0$ και $1$ (είναι λίγο μικρότερος του $1$).
\sol{18} Ελεύθερη απάντηση. Ο μαθητής πρέπει να παρατηρήσει ότι ο πιο αριστερός αριθμός στον άξονα είναι ο μικρότερος και να τοποθετήσει σωστά τα πρόσημα.
\sol{19} Ελεύθερη απάντηση. Π.χ. $8$ ώρες ύπνου: $\frac{8}{24}=\frac{1}{3}$ και $9$ ώρες: $\frac{9}{24}=\frac{3}{8}$. Για να συγκριθούν: $\frac{1}{3}=\frac{8}{24}<\frac{9}{24}=\frac{3}{8}$.
